\ncertchapter{15}{Rolling Motion}\chlabel{ch:rolling}

Rolling is where the two halves of mechanics meet. A rolling wheel translates and rotates at once, and the
two motions are not independent: they are locked together by a single constraint that holds as long as the
wheel does not slip. That constraint, $v_{\text{cm}}=\omega R$ (speed of the centre equals angular speed
times radius), is the subject of this chapter, and nearly everything follows from it: the split of kinetic
energy into translation and rotation, the direction of friction on a wheel, and a race down a ramp decided
by shape alone. First, though, it is worth seeing why rolling deserves a chapter of its own.

\begin{motivbox}{Why rolling matters: the wheel}
Rolling is the mechanics of the wheel, one of the most important machines ever built. Its value rests on one fact derived in this chapter: the contact point of a rolling wheel
is \emph{instantaneously at rest}, so its friction is \emph{static} and does \emph{no work}. Three
consequences follow:
\begin{itemize}
\item \textbf{Why the wheel beats the sledge.} Sliding a load of mass $m$ across level ground turns energy
$\mu_k mg$ into heat for every metre travelled (\chref{ch:friction}). Rolling it loses far less: only the
small losses of rolling resistance (discussed below) remain. That difference is why carts, trains, and
bearings exist.
\item \textbf{Why cars are driven by friction.} On a powered wheel the same static friction points
\emph{forward}, and it is what pushes the car along the road.
\item \textbf{Why shape decides a race.} Because a rolling body must divert energy into spinning, how its
mass is distributed, not how much it has nor how big it is, fixes how fast it rolls downhill: a marble
beats a hoop of any size.
\end{itemize}
\end{motivbox}

\section{The Rolling Constraint}

All three consequences rest on how the translation and the rotation of a wheel are tied together, so that is
where we begin. A wheel that rolls without slipping covers, in one turn, a distance equal to its
circumference. That observation is the whole of the rolling constraint. If the wheel, of radius $R$, turns
through an angle $\phi$ (in radians; we write $\phi$ rather than \chref{ch:rotation}'s $\theta$ because
$\theta$ will be the angle of an incline), the arc $R\phi$ of rim that has touched the ground equals the
distance $s$ the centre has moved: $s=R\phi$. Differentiating once and then again with respect to time gives
the velocity and acceleration forms below (valid for a round body rolling on a fixed, flat surface).

\begin{defbox}{Rolling without slipping}
When a body of radius $R$ rolls without its contact point sliding,
\[
v_{\text{cm}}=\omega R,\qquad a_{\text{cm}}=\alpha R,
\]
where $v_{\text{cm}}$ and $a_{\text{cm}}$ are the speed and acceleration of the centre of mass (CM), and
$\omega=d\phi/dt$ and $\alpha=d\omega/dt$ are the angular speed (\SI{}{rad/s}) and angular acceleration
(\SI{}{rad/s^2}) about it, as in \chref{ch:rotation}. $R$ is the radius of the part that touches the ground (for
a spool rolling on its inner axle, the axle's radius). This is a \emph{constraint}, not a law; it holds only
while the contact point does not slip.
\end{defbox}

\begin{derivbox}{Why the contact point is instantaneously at rest}
A rolling body's motion is a \emph{translation} of the whole body at $v_{\text{cm}}$ plus a
\emph{rotation} about its centre at $\omega$, in which a point on the rim moves at $\omega R$ tangent to the rim
(\chref{ch:circkin}). Add the two velocities at each point, as Figure~\ref{fig:rolladd} does arrow by arrow:
\stepn{1} \textbf{At the contact point} (bottom), translation gives $+v_{\text{cm}}$ (forward) and
rotation gives $\omega R$ \emph{backward}. If $v_{\text{cm}}=\omega R$ they cancel exactly and the
contact point is momentarily at rest. That cancellation \emph{is} the rolling condition.
\stepn{2} \textbf{At the top}, both point forward and add: the top of the wheel moves at
$v_{\text{cm}}+\omega R=2v_{\text{cm}}$, twice the speed of the axle. (This is why the top of a moving
car's tyre blurs in photographs while the bottom stays sharp.)
\stepn{3} \textbf{On a moving surface.} The real condition is that the body's contact point moves
\emph{with the surface it touches}. Taking forward as positive, that point moves at $v_{\text{cm}}-\omega R$,
so on a plank or belt moving at $v_{\text{s}}$ rolling without slipping means $v_{\text{cm}}-\omega R=v_{\text{s}}$;
fixed ground, $v_{\text{s}}=0$, gives back $v_{\text{cm}}=\omega R$.
\end{derivbox}

\begin{figure}[htbp]
\centering
\tikzfig{rolling_addition}
\figcaption{fig:rolladd}{Rolling is translation plus rotation. Adding the two velocity arrows at each point settles everything: at the contact point $P$ they are equal and opposite and cancel (that is $v_{\text{cm}}=\omega R$), while at the top they add to $2v_{\text{cm}}$. Every resultant is perpendicular to the dashed line from $P$ and proportional to its length: at each instant the wheel is pivoting about $P$.}
\end{figure}

\inwords{Because the contact point is momentarily at rest, you can picture the \emph{whole} rolling body,
at each instant, as \emph{pivoting} about it, as if hinged to the ground for a fleeting moment. The line
through the contact point, parallel to the axle, is the \textbf{instantaneous axis of rotation}: the speed
of any point is $\omega$ times its distance from it, so the top (distance $2R$) moves at $2v_{\text{cm}}$
and the centre (distance $R$) at $v_{\text{cm}}$. The axis moves along the ground as the wheel rolls. One
caution: the contact point has zero velocity but not zero acceleration (it accelerates upward at
$\omega^2R$), so the pivot picture is safe for velocities but not, in general, for accelerations.}

Knowing how fast every point moves, we can now ask how much energy the rolling body carries.

\begin{keybox}{Kinetic energy of a rolling body}
For a body of mass $M$ rolling without slipping, the total kinetic energy $K$ splits into translation of the
CM plus rotation about it (the split $K=\tfrac12Mv_{\text{cm}}^2+K'$ of \S\ref{sec:cmframe}, with $K'$
here the rotational energy $\tfrac12I_{\text{cm}}\omega^2$ of \chref{ch:rotation}):
\[
K=\tfrac12 Mv_{\text{cm}}^2+\tfrac12 I_{\text{cm}}\omega^2=\tfrac12 Mv_{\text{cm}}^2(1+c),
\]
using $\omega=v_{\text{cm}}/R$ and $I_{\text{cm}}=cMR^2$, the moment of inertia about the axle through the
CM written with \chref{ch:rotation}'s dimensionless shape factor $c$ (so the rotational term becomes
$\tfrac12 cMR^2\cdot v_{\text{cm}}^2/R^2=\tfrac12 cMv_{\text{cm}}^2$). The factor $(1+c)$ is the extra
``cost'' of having to spin as well as move: a solid sphere carries $c=\tfrac25$, a ring $c=1$.
(A second proof: about the instantaneous axis through the contact point $P$, the moment of inertia is
$I_P=I_{\text{cm}}+MR^2$ by the parallel-axis theorem, so
$K=\tfrac12 I_P\omega^2=\tfrac12 I_{\text{cm}}\omega^2+\tfrac12 Mv_{\text{cm}}^2$.)
\end{keybox}

\begin{trapbox}{the friction in pure rolling is static and does no work}
Is this kinetic energy drained as the body rolls? The only candidate is friction at the contact point.
That point is instantaneously at \emph{rest}, so the friction is \emph{static} (of size $f\le\mu_s N$, with
$N$ the normal force; usually an unknown to solve for) and, since its point of application does not move,
it does \emph{no work} on the body as a whole. It is not idle, though: on a slope,
over a distance $d$, it does work $-fd$ on the translation and $+fR\cdot(d/R)=+fd$ on the rotation, so it
moves energy from one form to the other without losing any. So mechanical energy is conserved in pure
rolling; only when a body \emph{skids} does kinetic friction act and dissipate energy.
\end{trapbox}

Friction may do no work, but it still exerts a force and a torque, so any force equation needs its
direction. Static friction takes whatever direction and size, up to $\mu_sN$, are needed to stop slipping
(\chref{ch:friction}). So which way does it point on a wheel?

\begin{methodbox}{The direction of static friction in rolling}
Friction opposes \emph{slipping}, not motion. Imagine the surface frictionless and ask which way the contact
point would then slide; friction points the other way.
\textbf{Free, on level ground at steady speed:} nothing tends to make it slip, so $f=0$; a rolling ball
needs no friction to keep rolling.
\textbf{Free, down a slope:} gravity alone would speed up the centre without spinning the body, so the
contact would slide downhill; friction acts \emph{up} the slope.
\textbf{Driven wheel:} the engine's torque would spin it faster than $v_{\text{cm}}/R$, so the contact would
slide backwards; friction acts \emph{forwards}. This forward friction is the only forward external force
on a car (as \chref{ch:com} required): the engine supplies the torque, the road does the pushing, and on
ice the wheels simply spin.
\textbf{Braked wheel:} the brake slows the spin, the contact would slide forwards, and friction acts backwards.
If you guess the direction wrongly, the algebra simply returns $f<0$.
\end{methodbox}

\begin{trapbox}{why real balls slow down (it is not sliding friction)}
If static friction does no work and is zero on level ground, why does a rolling ball eventually stop? The
cause is \textbf{rolling resistance}, a separate and much weaker effect. Real surfaces
\emph{deform} slightly at the contact (the tyre flattens, the ground dimples), and continually
re-forming that dimple as the wheel advances drains energy. The material ahead is being squashed while the
material behind does not spring back fully, so the ground pushes harder on the front of the contact patch:
the normal force acts slightly ahead of the centre and exerts a small torque that opposes the rolling. It
is why bicycle tyres are pumped hard (a stiffer tyre deforms less) and why a ball rolls further on marble
than on carpet.
\end{trapbox}

\section{Rolling Down an Incline}

With the friction's direction settled, we can put the pieces together on the standard rolling problem, a
body released on a slope. Solving it with forces and torques together gives
not only the acceleration but also the friction needed, and therefore the condition for rolling to occur at all.

\begin{derivbox}{Acceleration of a rolling body down an incline}
A body with $I_{\text{cm}}=cMR^2$ rolls without slipping down an incline of angle $\theta$; write $a$ for
$a_{\text{cm}}$, taken positive down the slope. The friction is static and of unknown size, so we treat it
as a second unknown and need two equations: one for translation, one for rotation.
\stepn{1} \textbf{Newton's second law along the incline.} The component $Mg\sin\theta$ of gravity acts down the slope and static friction $f$, assumed
\emph{up} the slope, opposes it (Step 2 confirms the sign, since it gives $f=cMa>0$):
$Mg\sin\theta-f=Ma$.
\stepn{2} \textbf{Torque about the CM.} Only friction has a moment arm about the centre (gravity and the
normal force act through it, as Figure~\ref{fig:rollforces} shows), so $fR=I_{\text{cm}}\alpha=cMR^2\cdot\dfrac{a}{R}$ (using $\alpha=a/R$).
This gives $f=cMa$. Taking torques about the CM is legitimate even though the CM accelerates: in the CM's
own frame every particle feels a pseudo-force $-m_i\va_{\text{cm}}$ (\S\ref{sec:noninertial}), and,
like uniform gravity, these act in effect at the CM and have no torque about it.
\stepn{3} \textbf{Combine.} Substitute $f=cMa$ into Step 1:
$Mg\sin\theta-cMa=Ma$, so
\[
\boxed{a=\frac{g\sin\theta}{1+c}.}
\]
\stepn{4} \textbf{Read the $(1+c)$.} It is a \emph{tax on rotation}. A sliding block ($c=0$) gets the full
$g\sin\theta$; a rolling body must divert some of the released energy into spinning up, so it accelerates less, and
the fraction of the kinetic energy diverted into spin is $c/(1+c)$, which is $\tfrac27$ for a solid
sphere and $\tfrac12$ for a ring.
\stepn{5} \textbf{Check about the contact point.} About $P$ only gravity has a moment arm, $R\sin\theta$,
so $MgR\sin\theta=I_P\alpha=(1+c)MR^2\cdot a/R$ (with $I_P=cMR^2+MR^2$ by the parallel-axis theorem):
the same $a$ in one line. (This works because, for a round
body, $P$ stays a fixed distance $R$ from the CM; about other moving points the short cut can fail.)
\end{derivbox}

\begin{figure}[htbp]
\centering
\tikzfig{rolling_forces}
\figcaption{fig:rollforces}{A body rolling down an incline, drawn to scale for a disc. Static friction acts \emph{up} the slope, and it is the only force with a moment arm about the centre, so it alone supplies the torque that spins the body up. The centre's acceleration $a=g\sin\theta/(1+c)$ down the slope is always less than the $g\sin\theta$ of a frictionless slide: the difference goes into rotation.}
\end{figure}

\begin{keybox}{The four results for a body rolling down an incline}
For any shape with $I_{\text{cm}}=cMR^2$ on an incline of angle $\theta$, everything follows from the
derivation above:
\[
a=\frac{g\sin\theta}{1+c},\qquad
v=\sqrt{\frac{2gh}{1+c}},\qquad
f=\frac{c}{1+c}\,Mg\sin\theta,\qquad
\mu_{\min}=\frac{c}{1+c}\tan\theta.
\]
The acceleration is the derivation's result; the speed $v$ after dropping a height $h$ follows from energy
conservation, $Mgh=\tfrac12Mv^2(1+c)$ (friction does no work); the friction $f=cMa$ is what the torque
equation requires; and $\mu_{\min}$ is the least coefficient of static friction
that can supply that $f$, from $f\le\mu_s Mg\cos\theta$ (below it, the body skids instead of rolling purely).
\emph{Every one is independent of $M$ and $R$}: only the shape $c$ enters.
\end{keybox}

\begin{figure}[htbp]
\centering
\tikzfig{rolling_race}
\figcaption{fig:rollrace}{The race is decided entirely by the shape factor $c$. Positions are shown at one instant after a common release, to scale (distance covered $\propto 1/(1+c)$). At any given height every body has the same kinetic energy per unit mass; what differs is how much of it goes into spinning. The small marble leads the large hoop, and a full tin of soup beats an empty one whatever their sizes.}
\end{figure}

Because these results depend on shape alone, releasing several bodies together gives a striking race;
Figure~\ref{fig:rollrace} catches it a moment after release.

\begin{exbox}{The race down the ramp}
A ring, a spherical shell, a disc, a solid sphere, and a block sliding without friction are released
together from the same height. Which arrives first? Since $a\propto\tfrac{1}{1+c}$ and
$v\propto\tfrac{1}{\sqrt{1+c}}$, everything hangs on $c$ (from \chref{ch:rotation}'s table):
\begin{center}
\begin{tabular}{@{}lccc@{}}
\toprule
Body & $c$ & $a/(g\sin\theta)=\tfrac{1}{1+c}$ & $v/\sqrt{2gh}=\tfrac{1}{\sqrt{1+c}}$\\
\midrule
Block (slides, no spin) & $0$ & $1.000$ & $1.000$\\
Solid sphere & $\tfrac25$ & $0.714$ & $0.845$\\
Disc / cylinder & $\tfrac12$ & $0.667$ & $0.816$\\
Spherical shell & $\tfrac23$ & $0.600$ & $0.775$\\
Ring / hoop & $1$ & $0.500$ & $0.707$\\
\bottomrule
\end{tabular}
\end{center}
So the order is \textbf{block $>$ solid sphere $>$ disc $>$ shell $>$ ring}, in order of increasing $c$.
The ring, with all its mass at radius $R$, has $c=1$: it splits its energy evenly between moving and
spinning, and its $a=\tfrac12 g\sin\theta$ is exactly half the frictionless value. Neither $M$ nor $R$
appears anywhere: a marble beats a hoop of any size, and ties with a boulder of the same shape.
\ans{smaller $c$ wins; independent of $M$ and $R$}
\end{exbox}

\why{Rolling is not a way of going \emph{faster}: the frictionless block wins the race. It is a way of
going \emph{cheaply}: the block wins only because we granted it zero friction, which no real surface
provides. Give it a coefficient of kinetic friction above $\tfrac12\tan\theta$ and it loses to every
roller (its acceleration $g(\sin\theta-\mu_k\cos\theta)$ falls below the ring's $\tfrac12 g\sin\theta$), because it turns energy into heat every metre while the rollers lose almost none.}

\section{Rolling with Slipping}\label{sec:slip15}

Everything so far assumed the contact does not slide. But pure rolling needs $\mu_s\ge\mu_{\min}$, and a
bowling ball launched without spin plainly slips at first. What happens then?

\begin{defbox}{Slip velocity}
The \textbf{slip velocity} is the velocity of the body's contact point over a fixed surface,
$v_{\text{cm}}-\omega R$ (forward positive). It is zero in pure rolling; positive when the body skids forward
($v_{\text{cm}}>\omega R$, a ball thrown without spin); negative when it spins too fast
($\omega R>v_{\text{cm}}$, a wheel spinning on ice).
\end{defbox}

While the body slips, the constraint $v_{\text{cm}}=\omega R$ is gone and the friction is \emph{kinetic},
of fixed size $\mu_kN$ and directed against the slip. Through its force and its torque it drives $v_{\text{cm}}$ and
$\omega R$ together, so on level ground the slip shrinks to zero; the friction then becomes static and, with
nothing left to oppose, vanishes, and the body rolls on purely. During the skid kinetic friction dissipates
energy, so energy is \emph{not} conserved; use the force and torque equations, or angular
momentum $L$ about a fixed point on the ground (Challenge~\ref{ex:skid15} works the bowling ball both ways).

On an incline with $\mu_s<\mu_{\min}$ rolling never becomes pure. Kinetic friction $\mu_kMg\cos\theta$ acts up
the slope, so Newton's second law gives $a=g(\sin\theta-\mu_k\cos\theta)$, while its torque,
$\mu_kMg\cos\theta\cdot R=cMR^2\alpha$, gives $\alpha=\mu_kg\cos\theta/(cR)$. Since
$\mu_k\le\mu_s<\tfrac{c}{1+c}\tan\theta$, a line of algebra shows $a>\alpha R$: the centre gains speed
faster than the rim, and the slip grows all the way down.

\section{Summary}\label{sec:summary15}
Rolling without slipping is the constraint $v_{\text{cm}}=\omega R$. Because the contact point is momentarily at
rest, the friction acting there is static and does no work, so mechanical energy is conserved for a body rolling
freely. Its kinetic energy splits into translation of the centre of mass and rotation about it, and the fraction
in rotation depends only on the body's shape.

\begin{itemize}
\item $K=\tfrac12Mv_{\text{cm}}^2(1+c)$ with $I_{\text{cm}}=cMR^2$: ring $c=1$, spherical shell $\tfrac23$, disc $\tfrac12$, solid sphere
      $\tfrac25$.
\item On an incline: $a=\dfrac{g\sin\theta}{1+c}$, $\;v=\sqrt{\dfrac{2gh}{1+c}}$, $\;f=\dfrac{c}{1+c}Mg\sin\theta$,
      $\;\mu_{\min}=\dfrac{c}{1+c}\tan\theta$.
\item The race down a ramp is decided by $c$ alone: frictionless block, then solid sphere, disc, shell, and ring.
\item On a freely rolling body friction points up the slope; on a driven wheel it points forward.
\item The top of a rolling wheel moves at $2v_{\text{cm}}$ and the contact point at zero.
\item Friction opposes slipping, not motion: zero for free rolling on level ground. While a body slips the
      friction is kinetic, dissipates energy, and drives $v_{\text{cm}}$ and $\omega R$ together until pure
      rolling begins.
\end{itemize}

\section{Problem Types}\label{sec:types15}
Five types cover rolling questions. Most reduce to the incline results above once the shape factor $c$ is
known, and the most efficient approach is usually energy rather than torque: the static friction at the contact does no work, so
$Mgh=\tfrac12Mv^2(1+c)$ gives the speed at once.

\archq{How the kinetic energy of a rolling body splits}
{``A ring, disc or sphere rolls without slipping; what fraction of its kinetic energy is rotational
(or translational)?''}
{Write $I=c MR^2$ ($c=1$ ring, $\tfrac12$ disc, $\tfrac25$ solid sphere, $\tfrac23$ shell).
Then $K_{\text{rot}}/K_{\text{trans}}=c$, so the rotational fraction of the kinetic energy is $\dfrac{c}{1+c}$ and the translational fraction $\dfrac{1}{1+c}$, independent of $M$, $R$ and speed.}
{Quoting the ratio $K_{\text{rot}}/K_{\text{trans}}=c$ when the question asks for the fraction of the \emph{total}.}
{For a solid sphere rolling without slipping, the fraction of its total kinetic energy that is
rotational is:}
{$\tfrac27$}{$\tfrac25$}{$\tfrac13$}{$\tfrac12$\archans{A}{With $c=\tfrac25$, the rotational
fraction is $\dfrac{c}{1+c}=\dfrac{2/5}{7/5}=\dfrac27$. (For a disc it is $\tfrac13$, for a ring
$\tfrac12$.)}}

\archq{Which body wins the race down an incline}
{``A ring, a disc, a solid sphere and a hollow sphere roll from the same height; which arrives first?'' or ``find the acceleration down an incline''.}
{Write $I=c MR^2$ and use $a=\dfrac{g\sin\theta}{1+c}$ and
$v=\sqrt{\dfrac{2gh}{1+c}}$. \textbf{Smaller $c$ wins}, and mass and radius cancel entirely:
sphere $(\tfrac25)$ $>$ disc $(\tfrac12)$ $>$ shell $(\tfrac23)$ $>$ ring $(1)$.}
{Thinking a heavier or bigger body wins. Only the \emph{shape}, through $c$, decides the order.}
{A ring rolls without slipping down a $30^\circ$ incline ($g=\SI{10}{m/s^2}$). Its acceleration is:}
{\SI{2.5}{m/s^2}}{\SI{3.3}{m/s^2}}{\SI{5}{m/s^2}}{\SI{10}{m/s^2}\archans{A}{For a ring $c=1$, so
$a=\dfrac{g\sin30^\circ}{1+1}=\dfrac{10(0.5)}{2}=\SI{2.5}{m/s^2}$, half the frictionless sliding
value.}}

\archq{Speed at the bottom from energy conservation}
{``A body rolls down from height $h$; find its speed at the bottom.''}
{Because friction does no work in \emph{pure} rolling, energy is conserved:
$Mgh=\tfrac12Mv^2(1+c)$, giving $v=\sqrt{\dfrac{2gh}{1+c}}$. Always \emph{less} than the sliding
value $\sqrt{2gh}$, because some energy goes into spin.}
{Using $v=\sqrt{2gh}$. That is the \emph{sliding} answer; a rolling body arrives slower.}
{A disc rolls without slipping from rest through a height of \SI{1.5}{m} ($g=\SI{10}{m/s^2}$). Its speed at the bottom is:}
{\SI{3.5}{m/s}}{\SI{4.5}{m/s}}{\SI{5.5}{m/s}}{\SI{7.7}{m/s}\archans{B}{
$v=\sqrt{\dfrac{2(10)(1.5)}{1.5}}=\sqrt{20}\approx\SI{4.5}{m/s}$. A sliding block would reach
$\sqrt{30}\approx\SI{5.5}{m/s}$.}}

\archq{Minimum friction needed for rolling without slipping}
{``What is the least $\mu_s$ that lets the body roll rather than slide?''}
{Required friction is $f=\dfrac{c Mg\sin\theta}{1+c}$, and $f\le\mu_s Mg\cos\theta$ gives
$\boxed{\mu_{\min}=\dfrac{c}{1+c}\tan\theta}$: $\tfrac27\tan\theta$ for a solid sphere,
$\tfrac13\tan\theta$ for a disc, $\tfrac12\tan\theta$ for a ring.}
{Assuming a smooth incline can support rolling. Without friction there is no torque about the centre, so
the body slides without ever starting to spin.}
{The minimum coefficient of friction for a solid sphere to roll down a $30^\circ$ incline is about:}
{$0.10$}{$0.16$}{$0.29$}{$0.58$\archans{B}{$\mu_{\min}=\tfrac27\tan30^\circ
=\tfrac27(0.577)\approx0.16$.}}

\archq{Velocity of a point on a rolling wheel}
{``Find the speed of the top/bottom/side point of a rolling wheel.''}
{Add the translation $v_{\text{cm}}$ and the rotation $\omega R$ \emph{as vectors} at that point. The
\textbf{contact point} is instantaneously at rest; the \textbf{top} moves at $2v_{\text{cm}}$; a point at the
side moves at $\sqrt2\,v_{\text{cm}}$.}
{Saying every point moves at $v_{\text{cm}}$. Only points a distance $R$ from the contact point (the centre
among them) do; the speed of any point is $\omega$ times its distance from the contact point.}
{For a wheel rolling without slipping at $v$, the speed of its topmost point is:}
{$0$}{$v$}{$\sqrt2\,v$}{$2v$\archans{D}{At the top, translation $(+v)$ and rotation $(+\omega R=v)$ point
the same way and add: $2v$. At the contact point they oppose and cancel to zero.}}

\section{Challenge Problems}\label{sec:challenge15}
The three problems below draw on all of Part D. Each contains a result that looks paradoxical at first
(energy lost and then regained; a collision that conserves angular momentum but not linear momentum; a final
speed independent of friction), and each is resolved by a careful application of a principle already
established.

\begin{challengebox}{The turntable and the energy from nowhere (builds on angular momentum)}
A disc of moment of inertia $I_0=\SI{4}{kg.m^2}$ spins freely at $\omega_0=\SI{3}{rad/s}$. A
\SI{2}{kg} lump of putty drops gently onto it at radius \SI{1}{m}. Find the new angular speed and the
energy lost. Then suppose the \SI{2}{kg} body is instead a toy that, after landing, crawls to the axis: find the final angular speed and the work the toy must do.
\challengeans{The putty lands vertically, exerting no torque about the spin axis, so \emph{angular
momentum is conserved}. The putty adds $I=mr^2=2(1)^2=\SI{2}{kg.m^2}$:
\[
I_0\omega_0=(I_0+mr^2)\omega_1\Rightarrow(4)(3)=(6)\omega_1\Rightarrow\omega_1=\SI{2}{rad/s}.
\]
Energy: $K_i=\tfrac12(4)(3)^2=\SI{18}{J}$, $K_f=\tfrac12(6)(2)^2=\SI{12}{J}$, so $\Delta K=\SI{-6}{J}$,
lost to friction as the putty skids up to speed. This is \chref{ch:collisions}'s perfectly inelastic collision in
angular clothing: $L$ conserved, $K$ dissipated.

\emph{Now the contrast.} Suppose the \SI{2}{kg} body is instead a small toy that, once moving with the disc at $\omega_1=\SI{2}{rad/s}$, crawls from $r=\SI{1}{m}$ to the axis. With no external torque, $L=\SI{12}{kg.m^2/s}$ is still conserved and $I$ falls back to \SI{4}{kg.m^2}, so $\omega_2=12/4=\SI{3}{rad/s}$ and $K=\tfrac12I\omega^2=L^2/2I=\SI{18}{J}$: the toy must do \SI{6}{J} of work, exactly the energy the landing dissipated (the same two values of $I$ appear in both steps). Both processes conserve $L$; one dissipates kinetic energy and the other gains it, because $L$ conservation says nothing about energy.}
\end{challengebox}

\begin{figure}[htbp]
\centering
\tikzfig{rolling_step}
\figcaption{fig:rollstep}{A rolling sphere at the instant it meets a step, drawn to scale for the challenge
($R=\SI{0.25}{m}$, $h=\SI{0.05}{m}$). The impulsive force acts at the edge $E$, so it has no torque about
$E$; the centre moves horizontally a perpendicular distance $R-h$ above $E$, which is the moment arm of
the sphere's linear momentum about $E$.}
\end{figure}

\begin{challengebox}{Rolling up a step (builds on angular momentum, impulse, rolling and energy)}
A uniform solid sphere ($c=\tfrac25$, mass $M$, radius $R$) rolls without slipping at speed $v_0$ along a
level floor and meets a step of height $h<R$ (Figure~\ref{fig:rollstep}). It strikes the edge of the step
and, without bouncing or slipping, begins to pivot about that edge. (a) Find its speed just after the impact. (b) Find the least
$v_0$ for which it climbs the step. (c) Show that it does not leave the edge at the start of the climb, and
find the fraction of kinetic energy lost in the impact. Evaluate for $R=\SI{0.25}{m}$, $h=\SI{0.05}{m}$,
$g=\SI{10}{m/s^2}$.
\challengeans{\textbf{(a) The impact.} The edge exerts a large impulsive force for a very short time, so
linear momentum is \emph{not} conserved. But that force acts \emph{at} the edge, so it has no torque about
the edge; gravity and the floor's normal force are finite and deliver negligible angular impulse during the
brief contact. Angular momentum about the edge is therefore conserved through the impact.

\emph{Before:} the CM moves horizontally at height $R$, a perpendicular distance $R-h$ above the edge,
giving orbital part $Mv_0(R-h)$; the spin part is $I_{\text{cm}}\omega_0=\tfrac25MR^2\cdot v_0/R$. So
$L=Mv_0\big(\tfrac75R-h\big)$. \emph{After:} the sphere turns about the edge with
$I_{\text{edge}}=\tfrac25MR^2+MR^2=\tfrac75MR^2$ (parallel axis) and CM speed $v_1=\omega_1R$, so
$L=\tfrac75MRv_1$. Equating,
\[
v_1=v_0\left(1-\frac{5h}{7R}\right).
\]
\textbf{(b) The climb.} Static friction at the edge does no work, so energy is conserved after the impact.
The CM rises by $h$ (from height $R-h$ above the edge to $R$), and all the kinetic energy is rotation about
the edge, $\tfrac12I_{\text{edge}}\omega_1^2=\tfrac12\cdot\tfrac75Mv_1^2$. Reaching the top needs
$\tfrac{7}{10}Mv_1^2\ge Mgh$, that is $v_1^2\ge\tfrac{10}{7}gh$, so
\[
v_{0,\min}=\frac{\sqrt{10gh/7}}{1-5h/7R}=\frac{7R\sqrt{10gh/7}}{7R-5h}.
\]
Numerically, $5h/7R=\tfrac17$, so $v_1=\tfrac67v_0$; $v_{1,\min}=\sqrt{5/7}=\SI{0.845}{m/s}$ and
$v_{0,\min}=\tfrac76(0.845)=\SI{0.986}{m/s}$.

\textbf{(c) Checks.} Just after impact the line from edge to centre makes an angle $\beta$ with the vertical,
$\cos\beta=(R-h)/R$. Along that line, towards the edge, the sphere's centre needs the centripetal force
$Mv_1^2/R$ of circular motion about the edge, supplied by gravity's component less the edge's normal force
$N$: $Mg\cos\beta-N=Mv_1^2/R$. Staying in contact needs $N\ge0$, that is
$v_1^2\le g(R-h)=\SI{2}{m^2/s^2}$, or $v_1\le\SI{1.41}{m/s}$. At the minimum, $v_1=\SI{0.845}{m/s}$ satisfies
this, and the condition only gets easier as the sphere rises (it slows and $\cos\beta$ grows), so the
sphere stays on the edge. (Faster than $v_1=\SI{1.41}{m/s}$, the sphere would hop off the edge, and the model
fails.) The kinetic energy after the impact is $\tfrac12\cdot\tfrac75Mv_1^2$ against $\tfrac12\cdot\tfrac75Mv_0^2$
before, so the fraction kept is $(1-5h/7R)^2=\tfrac{36}{49}$ and the fraction lost is
$\tfrac{13}{49}\approx27\%$, dissipated in the inelastic impact with the edge. Angular momentum about the
right point survives the collision; linear momentum and kinetic energy do not.}
\end{challengebox}

\begin{challengebox}{From skidding to rolling (builds on angular momentum and rolling)}\label{ex:skid15}
A bowling ball (a uniform solid sphere, $c=\tfrac25$, mass $M$, radius $R$) is launched along a level lane
with speed $v_0$ and no spin. The coefficient of kinetic friction between ball and lane is $\mu_k$. Find the
speed at which it starts to roll without slipping, the time this takes, and the fraction of its kinetic
energy lost on the way.
\challengeans{While the ball skids, its contact point slides forward over the lane, so kinetic friction
$f=\mu_k Mg$ acts \emph{backward}. It slows the translation, $v=v_0-\mu_k g t$, and its torque about the
centre, $fR=\tfrac25MR^2\alpha$, spins the ball up: $\omega=\dfrac{5\mu_k g}{2R}\,t$. Skidding stops when
$v=\omega R$:
\[
v_0-\mu_k g t=\tfrac52\mu_k g t\quad\Rightarrow\quad t=\frac{2v_0}{7\mu_k g},\qquad v=\tfrac57 v_0 .
\]
A quicker route avoids $\mu_k$ altogether. Take angular momentum about a fixed point on the lane, on the
line the contact point travels along. Friction acts along that line, so its torque about the point is zero;
gravity and the normal force are equal, opposite and act along the same vertical line, so their torques
cancel. Angular momentum about the point is therefore conserved (orbital part $MvR$ plus spin part
$\tfrac25MR^2\omega$):
\[
Mv_0R=MvR+\tfrac25MR^2\cdot\frac{v}{R}=\tfrac75MvR\quad\Rightarrow\quad v=\tfrac57v_0 .
\]
Energy: $K_f=\tfrac12Mv^2\big(1+\tfrac25\big)=\tfrac75\cdot\tfrac{25}{49}\cdot\tfrac12Mv_0^2=\tfrac57\cdot\tfrac12Mv_0^2$,
so $\tfrac27$ of the initial kinetic energy is lost as heat during the skid. Once rolling begins the
friction falls to zero (on level ground nothing tends to make the contact slip), and the ball rolls on at
$\tfrac57v_0$ until rolling resistance slowly stops it. The final speed does not depend on $\mu_k$: a
slippery lane only lengthens the skid.}
\end{challengebox}

\section*{Looking Back on Part D}
\addcontentsline{toc}{section}{Looking Back on Part D}

\textbf{Part D extended everything to bodies with size and spin.}
\begin{itemize}
\item \textbf{Centre of mass (\chref{ch:com}).} The CM is the mass-weighted average position, and it moves as a
single particle under the external force alone ($M\va_{\text{cm}}=\vF_{\text{ext}}$), because
internal forces cancel in third-law pairs. This is why a tumbling hammer's CM traces a clean parabola, and
the retroactive licence for every ``block'' drawn as a dot. With no external force the velocity of the CM cannot change: a person walking along a boat on still water pushes the boat back, and the CM of the pair stays put.
\item \textbf{Rotation (\chref{ch:rotation}).} Every translational law has a rotational twin: $x\to\theta$, $m\to I$,
$\vF\to\vtau=\vr\ex\vF$, $\vp\to\vL=I\vomega$. The one genuinely new object is $I=\sum mr^2$, and,
unlike mass, it depends on the axis. No external torque means $L=I\omega$ is fixed (the skater), and the
two axis theorems shift or combine $I$ without integrals.
\item \textbf{Rolling (\chref{ch:rolling}).} Translation plus rotation, locked by $v_{\text{cm}}=\omega R$: the contact
point is instantaneously at rest (so friction is static and does no work) and the top moves at $2v_{\text{cm}}$.
Down an incline, $a=g\sin\theta/(1+c)$, the $(1+c)$ a tax on rotation that decides the race by shape
alone. Friction opposes slipping, not motion; when a body does slip, kinetic friction dissipates energy until
pure rolling takes over.
\end{itemize}
\textbf{The thread into Part E.} We now have every tool: force, energy, momentum and angular momentum. The
final part aims them all at one force reaching across empty space.
